\documentclass[11pt]{article}
\usepackage[a4paper,margin=0.65in]{geometry}
\usepackage[T1]{fontenc}
\usepackage{amsmath,amssymb,amsfonts}
\usepackage{graphicx}
\usepackage{pdfpages}
\usepackage[english,french,shorthands=off]{babel}
\usepackage{fancyhdr}
\usepackage[bookmarks=false,colorlinks=false]{hyperref}
\usepackage[protrusion=true,expansion=false]{microtype}
\emergencystretch=3em
\tolerance=600
\providecommand{\modd}{\mathrm{mod}}
\providecommand{\Disc}{\mathrm{Disc}}
\providecommand{\canont}{}
\providecommand{\asterism}{*}
\providecommand{\Hessian}{H}
\providecommand{\tome}{}
\providecommand{\page}{p.}
\pagestyle{fancy}\fancyhf{}\fancyhead[L]{Pages 226--250}\fancyhead[R]{Cayley --- Collected Papers, Vol.~I}\renewcommand{\headrulewidth}{0.4pt}
\setlength{\headheight}{15pt}

\begin{document}

\begin{center}
{\LARGE\bfseries The Collected Mathematical Papers of}\\[0.5em]
{\LARGE\bfseries Arthur Cayley}\\[1em]
{\Large Volume I}\\[1em]
{\large PDF Pages 226--250 \quad (Book pages 204--228)}\\[2em]
\end{center}

\tableofcontents
\newpage

% ===========================================================================
% PDF page 226 -- Book page 204
% Paper [29] -- ADDITION A LA NOTE SUR QUELQUES INTEGRALES MULTIPLES.
% ===========================================================================
\section*{29.}
\section*{ADDITION \`A LA NOTE SUR QUELQUES INT\'EGRALES MULTIPLES.}
\addcontentsline{toc}{section}{29. Addition \`a la Note sur quelques Int\'egrales multiples}

\textit{[From the Journal de Math\'ematiques Pures et Appliqu\'ees (Liouville), tome x. (1845), pp.~242--244.]}

\bigskip
\selectlanguage{french}
\setcounter{page}{204}

On d\'emontre, au moyen des formules que j'ai donn\'ees sur ce sujet, [28], une propri\'et\'e remarquable de l'int\'egrale multiple
\[
V = \int dx_1 \ldots dx_n\, x_1^{h_1+1} \ldots x_n^{h_n+1} \phi(a_1 - x_1, \ldots),
\]
prise entre les limites
\[
\sum \frac{x_s}{a_s} = 1.
\]

En effet, en supposant toujours que l'on peut d\'evelopper la fonction $\phi$ selon les puissances enti\`eres et positives de $t$, l'int\'egrale $V$ peut s'exprimer sous la forme
\[
V = \tfrac{(-)^p t^{2p}}{2^{n-2}} \, S^{r=\infty}_{r=0}\!\left( \tfrac{d}{da_1}\right)^{\!h_1} \!\!\ldots \left( \tfrac{d}{da_n}\right)^{\!h_n}\!\!\cdot u \cdot t^{2r},
\]
o\`u
\[
k' = a_1 + \ldots + a_n,
\]
\[
u = \int_0^\infty \frac{x^{2v+1}}{2(\nu+1)\Gamma(p+k+f+\tfrac12 n+1)} \cdot \phi(a_1, \ldots).
\]
\[
V = \tfrac{(-)^p}{2^{n-2}} \cdot t^{2p} \cdot \frac{d^k}{da_1^{h_1}}\cdot
\]

Soit
\[
W = \int dx_1 \ldots\, \phi(a_1 - x_1 T, \ldots)
\]

\newpage
% ===========================================================================
% PDF page 227 -- Book page 205
% ===========================================================================
\setcounter{page}{205}

entre les m\^emes limites que $V$. On a
\[
W = \pi^{\tfrac12 n} t^{2p} \tfrac{1}{\Gamma(\tfrac12 n)} S^{r=\infty}_{r=0} \!\left( \tfrac{d}{da_1}\right)^{\!h_1} \!\!\ldots \left( \tfrac{d}{da_n}\right)^{\!h_n} T^{2p}\phi(a_1, \ldots).
\]

En multipliant par
\[
2T^{2p+1}(1-T^2)^{k+f}\,dT,
\]
et int\'egrant depuis $T = 0$ jusqu'\`a $T = 1$, on obtient
\[
2 \int_0^1 T^{2p+1}(1-T^2)^{k+f}\,W\,dT
\]
\[
= \pi^{\tfrac12 n} \tfrac{1}{\Gamma(\tfrac12 n)} S^{r=\infty}_{r=0} \!\left(\tfrac{d}{da_1}\right)^{\!h_1}\!\!\ldots\!\left(\tfrac{d}{da_n}\right)^{\!h_n}\!\!\cdot \frac{\Gamma(p+r+1)\Gamma(k+f+1)}{\Gamma(p+r+k+f+\tfrac12 n+1)}\,\phi(a_1, \ldots),
\]
puisque en g\'en\'eral
\[
\tfrac{\Gamma(k+f)}{\Gamma(\tfrac12 n)} \int_0^1 T^{2p+1}(1-T^2)^{k+f}\,dT = \tfrac{\Gamma(p+r+1)\Gamma(k+f+1)}{\Gamma(p+r+k+f+\tfrac12 n+1)}.
\]

On a donc l'\'equation
\[
\tfrac{\Gamma(k+f)}{\Gamma(\tfrac12 n)} \int_0^1 T^{2p+1}(1-T^2)^{k+f}\,W\,dT = \pi^{\tfrac12 n}\, t^{2p}\cdot U.
\]

Mettons la valeur de $U$ qui en r\'esulte, dans l'\'equation entre $V$ et $U$; faisons aussi $t = 1$, ce qui ne nuit pas \`a la g\'en\'eralit\'e. On a, en rassemblant les formules,
\[
V = \int dx_1 \ldots dx_n\, x_1^{h_1+1} \ldots x_n^{h_n+1} \phi(a_1 - x_1, \ldots),
\]
\[
W = \int dx_1 \ldots dx_n\, \phi(a_1 - x_1 T, \ldots),
\]
\[
V = \tfrac{(-)^p}{2^{n-2}} \tfrac{\Gamma(k+f)}{\Gamma(\tfrac12 n)}\!\left(\tfrac{d}{da_1}\right)^{\!h_1}\!\!\ldots\!\left(\tfrac{d}{da_n}\right)^{\!h_n}\!\!\cdot\!\int_0^1 T^{2p+1}(1-T^2)^{k+f}\,W\,dT,
\]
ce qui \'etablit ce th\'eor\`eme: La suite des int\'egrales $V$ s'exprime au moyen des coefficients diff\'erentiels par rapport \`a $A_1, \ldots A_n$, et \`a $a_1, \ldots, a_n$ de la seule int\'egrale
\[
\int_0^1 T^{2p+1}(1-T^2)^{k+f}\,dx_1 \ldots dx_n \ldots \phi(a_1 - x_1 T, \ldots).
\]

En supposant que les indices dans $V$ soient tous pairs, on a $f = 0$. Les symboles $\tfrac{d}{dA_1}, \ldots$ n'entrent plus dans l'expression de $V$. En changeant la fonction $\phi$, on a ces formules plus simples,

\newpage
% ===========================================================================
% PDF page 228 -- Book page 206
% ===========================================================================
\setcounter{page}{206}

\[
V = \int dx_1 \ldots dx_n\, x_1^{h_1} \ldots x_n^{h_n} \phi(x_1, \ldots, x_n),
\]
\[
W = \int dx_1 \ldots dx_n\, \phi(Tx_1, \ldots, Tx_n),
\]
\[
V = \tfrac{(-)^p}{2^{n-2}} \tfrac{\Gamma(k)}{\Gamma(\tfrac12 n)}\!\left(\tfrac{d}{da_1}\right)^{\!h_1}\!\!\ldots\!\left(\tfrac{d}{da_n}\right)^{\!h_n}\!\!\cdot\!\int_0^1 T^{2p+1}(1-T^2)^{k}\,W\,dT,
\]
o\`u $V$ s'exprime au moyen des coefficients diff\'erentiels par rapport \`a $A_1, \ldots A_n$ de l'int\'egrale
\[
\int_0^1 T^{2p+1}(1-T^2)^{k}\,dx_1\ldots dx_n\, \phi(Tx_1, \ldots, Tx_n),
\]
de mani\`ere que nous avons trouv\'e des relations entre des int\'egrales d\'efinies qui contiennent une fonction ind\'etermin\'ee, assujettie \`a la seule condition d'\^etre d\'eveloppable dans les limites de l'int\'egration selon les puissances enti\`eres et positives des variables\footnote{Cette condition est suppos\'ee dans la d\'emonstration, mais il me para\^it, d'apr\`es la forme du r\'esultat, qu'elle n'est pas essentielle.}; mais dans lesquelles les limites sont donn\'ees par
\[
\sum \frac{x_s}{a_s} = 1.
\]

\newpage
% ===========================================================================
% PDF page 229 -- Book page 207
% Paper [30] -- M\'EMOIRE SUR LES COURBES A DOUBLE COURBURE
% ===========================================================================
\section*{30.}
\section*{M\'EMOIRE SUR LES COURBES \`A DOUBLE COURBURE ET LES SURFACES D\'EVELOPPABLES.}
\addcontentsline{toc}{section}{30. M\'emoire sur les courbes \`a double courbure et les surfaces d\'eveloppables}

\textit{[From the Journal de Math\'ematiques Pures et Appliqu\'ees (Liouville), tome x. (1845), pp.~245--250.]}

\bigskip
\setcounter{page}{207}

On trouve, dans la \textit{Th\'eorie des courbes alg\'ebriques} de M.~Pl\"ucker, de tr\`es-belles recherches sur le nombre des diff\'erentes singularit\'es (points d'inflexion, tangentes doubles, \&c.) des courbes planes. Les m\^emes principes peuvent s'appliquer au cas des courbes \`a trois dimensions. Pour cela, consid\'erons une suite continue de points dans l'espace, les lignes qui passent par deux points cons\'ecutifs, et les plans qui passent par trois points cons\'ecutifs; ou, en envisageant autrement la m\^eme figure, une suite de lignes dont chacune rencontre la ligne cons\'ecutive, les points d'intersection de deux lignes cons\'ecutives, et les plans qui contiennent ces lignes; ou encore de cette mani\`ere: une suite de plans, les lignes d'intersection de deux plans cons\'ecutifs, les points d'intersection de trois plans cons\'ecutifs. C'est ce que l'on peut nommer un \textit{syst\`eme simple}. Ce syst\`eme est \'evidemment form\'e d'une courbe \`a double courbure et d'une surface d\'eveloppable; la courbe est l'ar\^ete de rebroussement de la surface, la surface est l'osculatrice d\'eveloppable de la courbe. Les points du syst\`eme sont des points dans la courbe, les lignes du syst\`eme sont les tangentes \`a la courbe, les plans du syst\`eme sont les plans osculateurs de la courbe. De m\^eme, les plans sont les plans tangents de la surface, les lignes sont les g\'en\'eratrices de la surface; pour les points, on peut les nommer les points de rebroussement de la surface. J'entendrai dans la suite par les termes \textit{ligne par deux points}, \textit{ligne dans deux plans}, une ligne men\'ee par deux points quelconques (non cons\'ecutifs en g\'en\'eral) du syst\`eme, et la ligne d'intersection de deux plans quelconques (non cons\'ecutifs en g\'en\'eral) du syst\`eme. Enfin, chaque ligne du syst\`eme est rencontr\'ee, en g\'en\'eral, par un certain nombre d'autres lignes non cons\'ecutives du syst\`eme. Je nommerai le point de rencontre d'une telle paire de lignes \textit{point dans deux lignes}, et le plan qui contient une telle paire \textit{plan par deux lignes}.

\newpage
% ===========================================================================
% PDF page 230 -- Book page 208
% ===========================================================================
\setcounter{page}{208}

Supposons qu'un plan donn\'e contient, en g\'en\'eral, $m$ points du syst\`eme, qu'une ligne donn\'ee rencontre, en g\'en\'eral, $r$ lignes du syst\`eme, qu'un point donn\'e est situ\'e, en g\'en\'eral, dans $n$ plans du syst\`eme. Le syst\`eme est dit \^etre de l'ordre $m$, du rang $r$, de la classe $n$. On voit tout de suite que l'ordre de la courbe est \'egal \`a l'ordre du syst\`eme, ou \`a $m$; la classe de la courbe au rang du syst\`eme, ou \`a $r$. Et de m\^eme, l'ordre de la surface au rang du syst\`eme, ou \`a $r$; la classe de la surface \`a la classe du syst\`eme, ou \`a $n$.

Cela pos\'e, les singularit\'es proprement dites (ouvrage cit\'e, page 202) sont les deux suivantes, analogues aux points d'inflexion et de rebroussement dans les courbes planes:

1. Quand quatre points cons\'ecutifs sont situ\'es dans le m\^eme plan, ou, autrement dit, quand trois lignes cons\'ecutives sont situ\'ees dans le m\^eme plan, ou quand deux plans cons\'ecutifs deviennent identiques; je dirai qu'il y a alors un \textit{plan stationnaire}, et je repr\'esenterai par la lettre $\alpha$ le nombre de ces plans.

2. Quand quatre plans cons\'ecutifs se rencontrent dans le m\^eme point, ou, autrement dit, quand trois lignes se rencontrent au m\^eme point, ou quand deux points cons\'ecutifs deviennent identiques; je dirai qu'il y a alors un \textit{point stationnaire}, et je repr\'esenterai par $\beta$ le nombre de ces points.

Dans le premier cas, il y a un point d'inflexion sph\'erique dans la courbe, et une ligne d'inflexion dans la surface. On n'a pas donn\'e de noms \`a ce qui arrive dans la courbe et la surface dans le second cas; et puisque la singularit\'e est suffisamment distingu\'ee d\'ej\`a en la nommant point stationnaire, il n'est pas n\'ecessaire de suppl\'eer \`a cette omission. On peut dire que ces deux cas sont les singularit\'es simples d'un syst\`eme. Il y a des singularit\'es d'un ordre plus \'elev\'e dont on n'a pas besoin ici. Ensuite, il y a des singularit\'es d'une autre esp\`ece, en quelque sorte analogues aux points et tangentes doubles, mais qui ont rapport \`a un point ou plan ind\'etermin\'e (hors du syst\`eme); savoir:

3. Un plan donn\'e peut contenir, \textit{en g\'en\'eral}, un nombre $g$ de \textit{lignes dans deux plans}.

4. Un point donn\'e peut \^etre situ\'e, \textit{en g\'en\'eral}, dans un nombre $h$ de \textit{lignes par deux points}.

5. Un plan donn\'e peut contenir, \textit{en g\'en\'eral}, un nombre $x$ de \textit{points dans deux lignes}.

6. Un point donn\'e peut \^etre situ\'e, \textit{en g\'en\'eral}, dans un nombre $y$ de \textit{plans par deux lignes}.

Ces quatre cas sont les singularit\'es impropres simples du syst\`eme.

Il faut maintenant chercher les relations qui ont lieu entre les nombres $m$, $r$, $n$, $\alpha$, $\beta$, $g$, $h$, $x$, $y$.

Citons d'abord les formules de M.~Pl\"ucker pour les courbes planes (page 211), en changeant seulement les lettres, pour \'eviter la confusion. En repr\'esentant par $\mu$ l'ordre d'une courbe, par $\nu$ sa classe, par $\xi$ le nombre de ses points doubles, $\eta$ de ses points

\newpage
% ===========================================================================
% PDF page 231 -- Book page 209
% ===========================================================================
\setcounter{page}{209}

de rebroussement, $\sigma$ de ses tangentes doubles, $\iota$ de ses points d'inflexion, l'on aura les six \'equations
\begin{align*}
\nu &= \mu\cdot\mu - 1 - (2\xi + 3\eta),\\
\iota &= 3\mu\cdot\mu - 2 - (6\xi + 8\eta),\\
\sigma &= \tfrac12 \mu\cdot\mu - 2 \mu^2 - 9 - (2\xi + 3\eta)(\mu\cdot\mu - 1 - 6) + 2\xi\cdot\xi - 1 + \tfrac32 \eta\cdot\eta - 1 + 6\xi\eta;\\
\mu &= \nu\cdot\nu - 1 - (2\sigma + 3\iota),\\
\eta &= 3\nu\cdot\nu - 2 - (6\sigma + 8\iota),\\
\xi &= \tfrac12 \nu\cdot\nu - 2 \nu^2 - 9 - (2\sigma + 3\iota)(\nu\cdot\nu - 1 - 6) + 2\sigma\cdot\sigma - 1 + \tfrac32 \iota\cdot\iota - 1 + 6\sigma\iota,
\end{align*}
dont les trois derni\`eres se d\'erivent des trois premi\`eres, et vice vers\^a.

Consid\'erons ensuite un plan donn\'e quelconque en conjonction avec le syst\`eme. Ce plan coupe la surface suivant une courbe plane. Les points de cette courbe sont les points de rencontre du plan avec les lignes du syst\`eme, les tangentes de cette courbe sont les lignes de rencontre du plan avec les plans du syst\`eme. Il est clair que la courbe est de l'ordre $r$ et de la classe $n$. Chaque fois que le plan contient un point en deux lignes, la courbe a un point double; aux points o\`u le plan rencontre la courbe \`a double courbure, il y a dans la courbe un rebroussement (car, dans ce cas, il y a deux lignes du syst\`eme qui coupent le plan en un m\^eme point, c'est-\`a-dire qu'il y a dans la courbe d'intersection un point stationnaire ou de rebroussement). Quand le plan contient une ligne en deux plans, il y a dans la courbe une tangente double; et enfin, pour chaque plan stationnaire du syst\`eme, il y a dans la courbe une tangente stationnaire ou une inflexion. Donc nous avons, dans la courbe, $r$ l'ordre, $n$ la classe, $x$ le nombre de points doubles, $m$ de points de rebroussement, $g$ de tangentes doubles, $\alpha$ de points d'inflexion, ce qui donne les six \'equations (\'equivalentes \`a trois):
\begin{align*}
n &= r\cdot r - 1 - (2x + 3m),\\
\alpha &= 3r\cdot r - 2 - (6x + 8m),\\
g &= \tfrac12 r\cdot r - 2r^2 - 9 - (2x + 3m)(r\cdot r - 1 - 6) + 2x\cdot x - 1 + \tfrac32 m\cdot m - 1 + 6xm;\\
r &= n\cdot n - 1 - (2g + 3\alpha),\\
m &= 3n\cdot n - 2 - (6g + 8\alpha),\\
x &= \tfrac12 n\cdot n - 2n^2 - 9 - (2g + 3\alpha)(n\cdot n - 1 - 6) + 2g\cdot g - 1 + \tfrac32 \alpha\cdot\alpha - 1 + 6g\alpha
\end{align*}
(o\`u l'on peut remarquer la sym\'etrie \`a l'\'egard des combinaisons $n, \alpha, g$, et $r, m, x$).

De m\^eme, en consid\'erant un point donn\'e quelconque en conjonction avec le syst\`eme, ce point d\'etermine, avec la courbe \`a double courbure, une surface conique, et, par des raisonnements pareils, on fait voir que, dans cette surface conique, l'ordre est $m$,

\newpage
% ===========================================================================
% PDF page 232 -- Book page 210
% ===========================================================================
\setcounter{page}{210}

la classe $r$, le nombre des lignes doubles $n$, des lignes de rebroussement $\beta$, des plans tangents doubles $y$, des lignes d'inflexion $h$, ce qui donne les \'equations (dont trois seulement sont ind\'ependantes):
\begin{align*}
r &= m\cdot m - 1 - (2h + 3\beta),\\
n &= 3m\cdot m - 2 - (6h + 8\beta),\\
y &= \tfrac12 m\cdot m - 2m^2 - 9 - (2h + 3\beta)(m\cdot m - 1 - 6) + 2h\cdot h - 1 + \tfrac32 \beta\cdot\beta - 1 + 6h\beta;\\
m &= r\cdot r - 1 - (2y + 3n),\\
\beta &= 3r\cdot r - 2 - (6y + 8n),\\
h &= \tfrac12 r\cdot r - 2r^2 - 9 - (2y + 3n)(r\cdot r - 1 - 6) + 2y\cdot y - 1 + \tfrac32 n\cdot n - 1 + 6yn
\end{align*}
(dans lesquelles on remarque la correspondance $r, n, y; m, \beta, h$: et, en les comparant avec les autres six \'equations, la correspondance $m, r, n, \alpha, \beta, g, h, x, y;\ n, r, m, \beta, \alpha, h, g, y, x$).

En consid\'erant une courbe \`a double courbure d'un ordre donn\'e $m$, on peut attribuer \`a $h, \beta$ des valeurs quelconques (entre certaines limites), et l'on a alors, pour d\'eterminer les autres quantit\'es, les \'equations suivantes:
\begin{align*}
r &= m\cdot m - 1 - (2h + 3\beta),\\
n &= 3m\cdot m - 2 - (6h + 8\beta),\\
y &= \tfrac12 m\cdot m - 2m^2 - 9 - (2h + 3\beta)(m\cdot m - 1 - 6) + 2h\cdot h - 1 + \tfrac32 \beta\cdot\beta - 1 + 6h\beta;\\
n &= r\cdot r - 1 - (2x + 3m),\\
\alpha &= 3r\cdot r - 2 - (6x + 8m),\\
g &= \tfrac12 r\cdot r - 2r^2 - 9 - (2x + 3m)(r\cdot r - 1 - 6) + 2x\cdot x - 1 + \tfrac32 m\cdot m - 1 + 6xm.
\end{align*}

Au cas d'une courbe plane les trois premi\`eres \'equations continuent d'\^etre vraies, mais les trois autres n'ont pas de sens. Le cas le plus simple est celui d'une courbe du troisi\`eme ordre dans l'espace. On a, dans ce cas, $m = 3$, $n = 1$, $\beta = 0$; et de l\`a le syst\`eme
\[
m, n, r, \alpha, \beta, g, h, x, y;
\]
\[
3, 3, 4, 0, 0, 1, 1, 0, 0;
\]
c'est-\`a-dire l'osculatrice d\'eveloppable d'une courbe du troisi\`eme ordre est une surface du quatri\`eme ordre, \&c. On obtient aussi un syst\`eme tr\`es-simple, en \'ecrivant $m = 4$, $h = 2$, $\beta = 1$, ce qui donne
\[
m, n, r, \alpha, \beta, g, h, x, y;
\]
\[
4, 4, 5, 1, 1, 2, 2, 2, 2.
\]

\newpage
% ===========================================================================
% PDF page 233 -- Book page 211
% ===========================================================================
\setcounter{page}{211}

Mais cela n'appartient pas, \`a ce que je crois, aux courbes les plus g\'en\'erales du quatri\`eme ordre.

Le probl\`eme de classifier les courbes \`a double courbure au moyen des surfaces que l'on peut faire passer par ces courbes, ou de trouver la nature d'une courbe qui est l'intersection de deux surfaces donn\'ees, para\^it appartenir plut\^ot \`a la th\'eorie des surfaces qu'\`a celle des courbes. Je n'ai rien de complet \`a offrir sur cela. Seulement je crois pouvoir dire que quand la courbe d'intersection de deux surfaces des ordres $\mu, \nu$ est de l'ordre $\mu\nu$ (ce qui est le cas g\'en\'eral), on a toujours
\[
2h = \mu\nu\cdot\overline{\mu - 1}\cdot\overline{\nu - 1},
\]
de mani\`ere que la classe de la courbe est au plus $mn(m+n-2)$. Mais j'esp\`ere revenir une autre fois sur cette question. Il est presque inutile de remarquer que pour un syst\`eme qui est r\'eciproque polaire d'un syst\`eme donn\'e, il faut seulement changer $m, r, n, \alpha, \beta, g, h, x, y$ en $n, r, m, \beta, \alpha, h, g, y, x$. Par exemple, les deux syst\`emes qui viennent d'\^etre consid\'er\'es ont des r\'eciproques de la m\^eme forme.

\newpage
% ===========================================================================
% PDF page 234 -- Book page 212
% Paper [31] -- DEMONSTRATION D'UN THEOREME DE M. CHASLES.
% ===========================================================================
\section*{31.}
\section*{D\'EMONSTRATION D'UN TH\'EOR\`EME DE M.~CHASLES.}
\addcontentsline{toc}{section}{31. D\'emonstration d'un Th\'eor\`eme de M. Chasles}

\textit{[From the Journal de Math\'ematiques Pures et Appliqu\'ees (Liouville), tome x. (1845), pp.~383--384.]}

\bigskip
\setcounter{page}{212}

``\,Soient $P, P'$ des points correspondants de deux figures homographiques; si la droite $PP'$ passe toujours par un point fixe $O$, les points $P$ sont situ\'es sur une courbe du troisi\`eme degr\'e, qui passe par ce m\^eme point.''

Soient $\frac{x}{w}, \frac{y}{w}, \frac{z}{w}$ les coordonn\'ees de $P$, $\frac{x'}{w'}, \frac{y'}{w'}, \frac{z'}{w'}$ celles de $P'$. En supposant que $x', y', z', w'$ sont des fonctions lin\'eaires (sans terme constant) de $x, y, z, w$, les deux figures seront homographiques.

Soient, de m\^eme, $\frac{\alpha}{\delta}, \frac{\beta}{\delta}, \frac{\gamma}{\delta}$ les coordonn\'ees de $O$; $\frac{\lambda}{\varpi}, \frac{\mu}{\varpi}, \frac{\nu}{\varpi}$ les coordonn\'ees d'un point quelconque $T$.

Puisque $P, P', O$ sont sur la m\^eme droite, on peut faire passer un plan par les quatre points $P, P', O, T$. Cela donne tout de suite l'\'equation
\[
\left| \begin{array}{cccc}
\lambda & \mu & \nu & \varpi\\
\alpha & \beta & \gamma & \delta\\
x' & y' & z' & w'\\
x & y & z & w
\end{array} \right| = 0
\]
(en repr\'esentant de cette mani\`ere le d\'eterminant form\'e avec les quantit\'es $\lambda, \mu, \nu$, \&c., $x', y', z', w'$),
\'equation qui doit \^etre satisfaite quels que soient $\lambda, \mu, \nu, \varpi$, et qui \'equivaut ainsi aux deux conditions
\[
\left| \begin{array}{ccc}
\alpha & \beta & \gamma\\
x' & y' & z'\\
x & y & z
\end{array} \right| = 0, \qquad
\left| \begin{array}{ccc}
\alpha & \beta & \delta\\
x' & y' & w'\\
x & y & w
\end{array} \right| = 0.
\]

Ces deux derni\`eres \'equations sont du second degr\'e par rapport aux quantit\'es $\frac{x}{w}, \frac{y}{w}, \frac{z}{w}$. Le point $P$ est donc situ\'e \`a l'intersection de deux surfaces du second ordre. Mais ces surfaces ont en commun la droite repr\'esent\'ee par les \'equations
\[
\alpha y - \beta x = 0, \qquad \alpha w - \delta x = 0.
\]
Donc elles se coupent de plus suivant une courbe du troisi\`eme degr\'e qui passe \'evidemment par le point $O$, parce qu'on satisfait aux \'equations en \'ecrivant
\[
x : \alpha = y : \beta = z : \gamma = w : \delta.
\]

\selectlanguage{english}

\newpage
% ===========================================================================
% PDF page 235 -- Book page 213
% Paper [32] -- ON SOME ANALYTICAL FORMULAE
% ===========================================================================
\section*{32.}
\section*{ON SOME ANALYTICAL FORMUL\AE, AND THEIR APPLICATION TO THE THEORY OF SPHERICAL COORDINATES.}
\addcontentsline{toc}{section}{32. On some Analytical Formul\ae{}, and their application to the Theory of Spherical Coordinates}

\textit{[From the Cambridge and Dublin Mathematical Journal, vol.~i. (1846), pp.~22--33.]}

\bigskip
\setcounter{page}{213}

\subsection*{Section 1.}

The formul\ae{} in question are only very particular cases of some relating to the theory of the transformation of functions of the second order, which will be given in a following paper. But the case of three variables, here as elsewhere, admits of a symmetrical notation so much simpler than in any other case (on the principle that with three quantities $a, b, c$, functions of $b, c$; of $c, a$; and of $a, b$, may symmetrically be denoted by $A, B, C$, which is not possible with a greater number of variables) that it will be convenient to employ here a notation entirely different from that made use of in the general case, and by means of which the results will be exhibited in a more compact form. There is no difficulty in verifying by actual multiplication, any of the equations here obtained.

It will be expedient to employ the abbreviation of making a single letter stand for a system of quantities. Thus for instance, if $\Phi = \theta, \phi, \psi$, this merely means that $\Phi(\xi)$ is to stand for $\Phi(\theta, \phi, \psi)$, $k\Phi$ for $k\theta, k\phi, k\psi$, \&c.

Suppose then
\[
\omega = \xi, \eta, \zeta \qquad (1),
\]
\[
\omega' = \xi', \eta', \zeta',
\]
\[
Q = A, B, C, F, G, H \qquad (2),
\]
\[
W(\omega, \omega', Q) = A\xi\xi' + B\eta\eta' + C\zeta\zeta' + F(\eta\zeta' + \eta'\zeta) + G(\zeta\xi' + \zeta'\xi) + H(\xi\eta' + \xi'\eta) \qquad (3),
\]
the function $W$ satisfies a remarkable equation, as follows:

\newpage
% ===========================================================================
% PDF page 236 -- Book page 214
% ===========================================================================
\setcounter{page}{214}

write
\begin{align*}
\mathfrak{A} &= BC - F^2, \qquad (4),\\
\mathfrak{B} &= CA - G^2,\\
\mathfrak{C} &= AB - H^2,\\
\mathfrak{F} &= GH - AF,\\
\mathfrak{G} &= HF - BG,\\
\mathfrak{H} &= FG - CH.
\end{align*}
\[
\mathbf{Q} = \mathfrak{A}, \mathfrak{B}, \mathfrak{C}, \mathfrak{F}, \mathfrak{G}, \mathfrak{H} \qquad (5),
\]
\[
\varpi = \eta\zeta' - \eta'\zeta,\ \zeta\xi' - \zeta'\xi,\ \xi\eta' - \xi'\eta \qquad (6),
\]
we have
\[
W(\omega_1, \omega_2, Q)\,W(\omega_3, \omega_4, Q) - W(\omega_1, \omega_3, Q)\,W(\omega_2, \omega_4, Q) = W(\varpi_1, \varpi_2, \mathbf{Q}) \qquad (7);
\]
of which we may notice also the particular cases
\[
W(\omega_1, \omega_2, Q)\,W(\omega_3, \omega_3, Q) - W(\omega_1, \omega_3, Q)\,W(\omega_2, \omega_3, Q) = W(\varpi_1, \varpi_2, \mathbf{Q}) \qquad (8),
\]
\[
W(\omega_1, \omega_1, Q)\,W(\omega_2, \omega_2, Q) - \{W(\omega_1, \omega_2, Q)\}^2 = W(\varpi_1, \varpi_2, \mathbf{Q}) \qquad (9).
\]

To these we may join the following formul\ae{}, for the transformation of the function $W$. Suppose
\[
\omega_1 = ax_1 + a'y_1 + a''z_1,\ bx_1 + b'y_1 + b''z_1,\ cx_1 + c'y_1 + c''z_1 \qquad (10),
\]
\[
\omega_2 = ax_2 + a'y_2 + a''z_2,\ bx_2 + b'y_2 + b''z_2,\ cx_2 + c'y_2 + c''z_2,
\]
then, writing
\[
g = a, b, c \qquad (11),
\]
\[
g' = a', b', c',
\]
\[
g'' = a'', b'', c'',
\]
\[
p_1 = x_1, y_1, z_1 \qquad (12),
\]
\[
p_2 = x_2, y_2, z_2,
\]
\[
\Theta = W(g, g, Q), W(g', g', Q), W(g'', g'', Q), W(g', g'', Q), W(g'', g, Q), W(g, g', Q) \qquad (13),
\]
we have
\[
W(\omega_1, \omega_2, Q) = W(p_1, p_2, \Theta) \qquad (14).
\]

Similarly, writing
\[
\boldsymbol{\Theta} = W(\varpi', \varpi'', \mathbf{Q}), W(\varpi'', \varpi, \mathbf{Q}), W(\varpi, \varpi', \mathbf{Q}),
\]
\[
W(\varpi'', \varpi, \mathbf{Q}), W(\varpi, \varpi', \mathbf{Q}), W(\varpi', \varpi'', \mathbf{Q}) \qquad (15),
\]
we have
\[
W(\varpi_1, \varpi_2, \mathbf{Q}) = W(p_1, p_2, \boldsymbol{\Theta}) \qquad (16),
\]
in which equations $\mathbf{Q}$ may obviously be changed into $Q$.

\newpage
% ===========================================================================
% PDF page 237 -- Book page 215
% ===========================================================================
\setcounter{page}{215}

\subsection*{Section 2. Geometrical Applications.}

Consider any three axes $Ax, Ay, Az$, and let $\lambda, \mu, \nu$ be the cosines of the inclinations of these lines to each other.

Let $\Lambda, M, N$ be the inclinations of the coordinate planes to each other; $l, m, n$ the inclination of the axes to the coordinate planes. Suppose, besides,
\begin{align*}
\mathfrak{a} &= 1 - \lambda^2,\\
\mathfrak{b} &= 1 - \mu^2,\\
\mathfrak{c} &= 1 - \nu^2,\\
\mathfrak{f} &= \mu\nu - \lambda,\\
\mathfrak{g} &= \nu\lambda - \mu,\\
\mathfrak{h} &= \lambda\mu - \nu, \qquad (18);\\
k &= 1 - \lambda^2 - \mu^2 - \nu^2 + 2\lambda\mu\nu
\end{align*}
we have the following systems of equations:
\begin{align*}
\sqrt{(\mathfrak{b}\mathfrak{c})}\cos\Lambda &= -\mathfrak{f}, & \sqrt{(\mathfrak{b}\mathfrak{c})}\sin\Lambda &= \sqrt{(k)}, & \sqrt{(\mathfrak{a})}\sin l &= \sqrt{(k)} \qquad (19),\\
\sqrt{(\mathfrak{c}\mathfrak{a})}\cos M &= -\mathfrak{g}, & \sqrt{(\mathfrak{c}\mathfrak{a})}\sin M &= \sqrt{(k)}, & \sqrt{(\mathfrak{b})}\sin m &= \sqrt{(k)}\\
\sqrt{(\mathfrak{a}\mathfrak{b})}\cos N &= -\mathfrak{h}, & \sqrt{(\mathfrak{a}\mathfrak{b})}\sin N &= \sqrt{(k)}, & \sqrt{(\mathfrak{c})}\sin n &= \sqrt{(k)}.
\end{align*}
\begin{align*}
\mathfrak{a} + \nu\mathfrak{h} + \mu\mathfrak{g} &= k, \qquad (20),\\
\nu\mathfrak{a} + \mathfrak{h} + \lambda\mathfrak{g} &= 0,\\
\mu\mathfrak{a} + \lambda\mathfrak{h} + \mathfrak{g} &= 0.
\end{align*}
\begin{align*}
\mathfrak{h} + \nu\mathfrak{b} + \mu\mathfrak{f} &= 0, \qquad (21),\\
\nu\mathfrak{h} + \mathfrak{b} + \lambda\mathfrak{f} &= k,\\
\mu\mathfrak{h} + \lambda\mathfrak{b} + \mathfrak{f} &= 0.
\end{align*}
\begin{align*}
\mathfrak{g} + \nu\mathfrak{f} + \mu\mathfrak{c} &= 0, \qquad (22),\\
\nu\mathfrak{g} + \mathfrak{f} + \lambda\mathfrak{c} &= 0,\\
\mu\mathfrak{g} + \lambda\mathfrak{f} + \mathfrak{c} &= k.
\end{align*}
\begin{align*}
\mathfrak{b}\mathfrak{c} - \mathfrak{f}^2 &= k\mathfrak{a}, \qquad (23),\\
\mathfrak{c}\mathfrak{a} - \mathfrak{g}^2 &= k\mathfrak{b},\\
\mathfrak{a}\mathfrak{b} - \mathfrak{h}^2 &= k\mathfrak{c},\\
\mathfrak{g}\mathfrak{h} - \mathfrak{a}\mathfrak{f} &= k\mathfrak{f},\\
\mathfrak{h}\mathfrak{f} - \mathfrak{b}\mathfrak{g} &= k\mathfrak{g},\\
\mathfrak{f}\mathfrak{g} - \mathfrak{c}\mathfrak{h} &= k\mathfrak{h},
\end{align*}
\[
\mathfrak{a}\mathfrak{b}\mathfrak{c} - \mathfrak{a}\mathfrak{f}^2 - \mathfrak{b}\mathfrak{g}^2 - \mathfrak{c}\mathfrak{h}^2 + 2\mathfrak{f}\mathfrak{g}\mathfrak{h} = k^2 \qquad (24).
\]

\newpage
% ===========================================================================
% PDF page 238 -- Book page 216
% ===========================================================================
\setcounter{page}{216}

Imagine now a line $AO$, and let $\alpha, \beta, \gamma$ be the cosines of its inclinations to the three axes. Suppose also, $\theta, \phi, \chi$ being its inclinations to the coordinate planes, we write
\[
\mathfrak{a} = \frac{\sin\theta}{\sqrt{(\mathfrak{a})}}, \quad \mathfrak{b} = \frac{\sin\phi}{\sqrt{(\mathfrak{b})}}, \quad \mathfrak{c} = \frac{\sin\chi}{\sqrt{(\mathfrak{c})}} \qquad (25).
\]

If we consider a point $P$ on the line $AO$, at a distance unity from the origin, we see immediately, by considering the projections in the directions perpendicular to the coordinate planes, that the coordinates of this point are $\mathfrak{a}, \mathfrak{b}, \mathfrak{c}$. By projecting on the three axes and on the line $AO$, we then obtain the equations
\begin{align*}
\alpha &= \mathfrak{a} + \nu\mathfrak{b} + \mu\mathfrak{c} \qquad (26),\\
\beta &= \nu\mathfrak{a} + \mathfrak{b} + \lambda\mathfrak{c},\\
\gamma &= \mu\mathfrak{a} + \lambda\mathfrak{b} + \mathfrak{c},\\
1 &= \alpha\mathfrak{a} + \beta\mathfrak{b} + \gamma\mathfrak{c} \qquad (27),
\end{align*}
from which we obtain
\begin{align*}
k\mathfrak{a} &= \mathfrak{a}\alpha + \mathfrak{h}\beta + \mathfrak{g}\gamma \qquad (28),\\
k\mathfrak{b} &= \mathfrak{h}\alpha + \mathfrak{b}\beta + \mathfrak{f}\gamma,\\
k\mathfrak{c} &= \mathfrak{g}\alpha + \mathfrak{f}\beta + \mathfrak{c}\gamma,\\
1 &= \alpha\mathfrak{a} + \beta\mathfrak{b} + \gamma\mathfrak{c} \qquad (29),
\end{align*}
and hence
\[
1 = \mathfrak{a}^2 + \mathfrak{b}^2 + \mathfrak{c}^2 + 2\lambda\mathfrak{b}\mathfrak{c} + 2\mu\mathfrak{a}\mathfrak{c} + 2\nu\mathfrak{a}\mathfrak{b} \qquad (30),
\]
\[
k = \mathfrak{a}\alpha^2 + \mathfrak{b}\beta^2 + \mathfrak{c}\gamma^2 + 2\mathfrak{f}\beta\gamma + 2\mathfrak{g}\alpha\gamma + 2\mathfrak{h}\alpha\beta \qquad (31).
\]

Hence writing
\begin{align*}
\mathfrak{a}, \mathfrak{b}, \mathfrak{c} &= t \qquad (32),\\
\alpha, \beta, \gamma &= \tau \qquad (33),\\
1, 1, 1, \lambda, \mu, \nu &= q \qquad (34),\\
\mathfrak{a}, \mathfrak{b}, \mathfrak{c}, \mathfrak{f}, \mathfrak{g}, \mathfrak{h} &= \mathbf{q} \qquad (35),
\end{align*}
we have the equations
\[
1 = W(t, t, q) \qquad (36),
\]
\[
k = W(\tau, \tau, \mathbf{q}) \qquad (37).
\]
Let $AO'$ be any other line, and $\delta$ its inclination to $AO$: $\alpha', \beta', \gamma', \mathfrak{a}', \mathfrak{b}', \mathfrak{c}'$, the quantities corresponding to $\alpha, \beta, \gamma, \mathfrak{a}, \mathfrak{b}, \mathfrak{c}$, and similarly $t', \tau'$ to $t, \tau$. We have of course
\[
1 = W(t', t', q) \qquad (38),
\]
\[
k = W(\tau', \tau', \mathbf{q}) \qquad (39).
\]

\newpage
% ===========================================================================
% PDF page 239 -- Book page 217
% ===========================================================================
\setcounter{page}{217}

We have besides, by projecting on the line $AO'$, the equation
\[
\cos\delta = \alpha'\mathfrak{a} + \beta'\mathfrak{b} + \gamma'\mathfrak{c} \qquad (40),
\]
or the analogous one
\[
\cos\delta = \alpha\mathfrak{a}' + \beta\mathfrak{b}' + \gamma\mathfrak{c}' \qquad (41).
\]
From either of which we deduce
\[
\cos\delta = \mathfrak{a}\mathfrak{a}' + \mathfrak{b}\mathfrak{b}' + \mathfrak{c}\mathfrak{c}' + \lambda(\mathfrak{b}\mathfrak{c}' + \mathfrak{b}'\mathfrak{c}) + \mu(\mathfrak{c}\mathfrak{a}' + \mathfrak{c}'\mathfrak{a}) + \nu(\mathfrak{a}\mathfrak{b}' + \mathfrak{a}'\mathfrak{b}) \qquad (42),
\]
\[
k\cos\delta = \mathfrak{a}\alpha\alpha' + \mathfrak{b}\beta\beta' + \mathfrak{c}\gamma\gamma' + \mathfrak{f}(\beta\gamma' + \beta'\gamma) + \mathfrak{g}(\gamma\alpha' + \gamma'\alpha) + \mathfrak{h}(\alpha\beta' + \alpha'\beta) \qquad (43);
\]
which may otherwise be written
\[
\cos\delta = W(t, t', q) \qquad (44),
\]
\[
k\cos\delta = W(\tau, \tau', \mathbf{q}) \qquad (45),
\]
or again, observing the equations which connect the quantities $t, \tau$,
\[
\cos\delta = \frac{W(t, t', q)}{\sqrt{\{W(t, t, q)\,W(t', t', q)\}}} \qquad (46),
\]
\[
\cos\delta = \frac{W(\tau, \tau', \mathbf{q})}{\sqrt{\{W(\tau, \tau, \mathbf{q})\,W(\tau', \tau', \mathbf{q})\}}} \qquad (47);
\]
forms which, though more complicated, have certain advantages; for instance, we derive immediately from them the new equations
\[
\sin\delta = \frac{\sqrt{\{W(\varpi, \varpi, q)\}}}{\sqrt{\{W(t, t, q)\,W(t', t', q)\}}} \qquad (48),
\]
\[
\sin\delta = \frac{\sqrt{\{W(\varpi', \varpi', \mathbf{q})\}}}{\sqrt{\{W(\tau, \tau, \mathbf{q})\,W(\tau', \tau', \mathbf{q})\}}}
\]
written more simply thus
\[
\sin\delta = \sqrt{\{W(\varpi, \varpi, q)\}} \qquad (49),
\]
\[
\sqrt{k}\sin\delta = \sqrt{\{W(\varpi', \varpi', \mathbf{q})\}} \qquad (50);
\]
to these we may join
\[
\cot\delta = \frac{W(t, t', q)}{\sqrt{\{W(\varpi, \varpi, q)\}}} \qquad (52),
\]
\[
\sqrt{k}\cot\delta = \frac{W(\tau, \tau', \mathbf{q})}{\sqrt{\{W(\varpi', \varpi', \mathbf{q})\}}} \qquad (53).
\]

\subsection*{Section 3. On Spherical Coordinates.}

Consider the points $X, Y, Z$, on the surface of a sphere, as the intersections of the three axes of the preceding section, with a sphere having its centre in the origin. It is evident that $\lambda, \mu, \nu$ are the cosines of the sides of the spherical triangle $XYZ$,

\newpage
% ===========================================================================
% PDF page 240 -- Book page 218
% ===========================================================================
\setcounter{page}{218}

$\Lambda, M, N$ are its sides, $l, m, n$ are the perpendiculars from the angles upon the opposite sides. Let $P$ be the point where the line $AO$ intersects the sphere: the position of the point $P$ may be determined by means of the ratios $\xi : \eta : \zeta$, supposing $\xi, \eta, \zeta$ denote quantities proportional to the $\alpha, \beta, \gamma$ of the preceding section, i.e.
\[
\xi : \eta : \zeta = \cos PX : \cos PY : \cos PZ \qquad (54);
\]
or again, by means of the ratios $x : y : z$, supposing $x, y, z$ denote quantities proportional to the $\mathfrak{a}, \mathfrak{b}, \mathfrak{c}$ of the preceding section, i.e.
\[
x : y : z = \frac{\sin Px}{\sin X} : \frac{\sin Py}{\sin Y} : \frac{\sin Pz}{\sin Z} \qquad (55),
\]
($Px, Py, Pz$ are the perpendiculars from $P$ on the sides of the spherical triangle $XYZ$).

These last equations may be otherwise written,
\[
\frac{x\sin X}{y\sin Y} = \frac{\sin PZY}{\sin PZX} \qquad (56),
\]
\[
\frac{y\sin Y}{z\sin Z} = \frac{\sin PXZ}{\sin PXY},
\]
\[
\frac{z\sin Z}{x\sin X} = \frac{\sin PYX}{\sin PYZ}.
\]

The ratios $\xi : \eta : \zeta$, or $x : y : z$, are termed the spherical coordinate ratios of the point $P$. The two together may be termed conjoint systems: the first may be termed the cosine system, and the second the sine system. The coordinates of the two systems are evidently connected by
\[
\xi : \eta : \zeta = x + \nu y + \mu z : \nu x + y + \lambda z : \mu x + \lambda y + z \qquad (57),
\]
or
\[
x : y : z = \mathfrak{a}\xi + \mathfrak{h}\eta + \mathfrak{g}\zeta : \mathfrak{h}\xi + \mathfrak{b}\eta + \mathfrak{f}\zeta : \mathfrak{g}\xi + \mathfrak{f}\eta + \mathfrak{c}\zeta \qquad (58).
\]
The systems may conveniently be represented by the single letters
\[
\omega = \xi, \eta, \zeta \qquad (59),
\]
\[
p = x, y, z \qquad (60).
\]

\subsection*{Fundamental formul\ae{} of spherical coordinates; distance of two points.}

Let $P, P'$ be the points, $\delta$ their distance, $\omega, p$ the conjoint coordinate systems of the first point, $\omega', p'$ of the second; we have obviously
\[
\cos\delta = \frac{W(p, p', q)}{\sqrt{\{W(p, p, q)\,W(p', p', q)\}}} \qquad (61),
\]
\[
\sqrt{k}\sin\delta = \frac{\sqrt{\{W(pp', pp', \mathbf{q})\}}}{\sqrt{\{W(p, p, q)\,W(p', p', q)\}}},
\]
\[
\cot\delta = \frac{W(p, p', q)}{\sqrt{\{W(pp', pp', \mathbf{q})\}}},
\]

\newpage
% ===========================================================================
% PDF page 241 -- Book page 219
% ===========================================================================
\setcounter{page}{219}

or
\[
\cos\delta = \frac{W(\omega, \omega', \mathbf{q})}{\sqrt{\{W(\omega, \omega, \mathbf{q})\,W(\omega', \omega', \mathbf{q})\}}} \qquad (62),
\]
\[
\frac{1}{\sqrt{k}}\sin\delta = \frac{\sqrt{\{W(\varpi\varpi', \varpi\varpi', q)\}}}{\sqrt{\{W(\omega, \omega, \mathbf{q})\,W(\omega', \omega', \mathbf{q})\}}},
\]
\[
\sqrt{k}\cot\delta = \frac{W(\omega, \omega', \mathbf{q})}{\sqrt{\{W(\varpi\varpi', \varpi\varpi', q)\}}}.
\]

\subsection*{Equation of a great Circle.}

Let the conjoint coordinate systems of the pole be
\[
e = a, b, c \qquad (63),
\]
\[
\mathbf{e} = \alpha, \beta, \gamma \qquad (64),
\]
then, expressing that the distance of any point $P$ in the locus from the pole is equal to $90^\circ$, we have immediately the equations
\[
W(p, \mathbf{e}, q) = 0 \qquad (65),
\]
\[
W(\omega, e, q) = 0 \qquad (66),
\]
which may otherwise be written in the forms
\[
\alpha\xi + \beta\eta + \gamma\zeta = 0 \qquad (67),
\]
\[
ax + by + cz = 0 \qquad (68),
\]
or the equation of a great circle is linear in either coordinate system. Conversely, any linear equation belongs to a great circle.

Suppose the equation given in the form
\[
A\xi + B\eta + C\zeta = 0 \qquad (69);
\]
or by an equation between cosine coordinate ratios:--- the sine system for the pole is given by
\[
\mathbf{e} = A, B, C \qquad (70),
\]
and the cosine system by
\[
e = A + \nu B + \mu C,\ \nu A + B + \lambda C,\ \mu A + \lambda B + C \qquad (71).
\]
Suppose the circle given by an equation between sine coordinates, or in the form
\[
Ax + By + Cz = 0 \qquad (72),
\]
the cosine system of coordinates for the pole is given by
\[
\mathbf{e} = A, B, C \qquad (73),
\]
and the sine system by
\[
e = \mathfrak{a}A + \mathfrak{h}B + \mathfrak{g}C,\ \mathfrak{h}A + \mathfrak{b}B + \mathfrak{f}C,\ \mathfrak{g}A + \mathfrak{f}B + \mathfrak{c}C \qquad (74).
\]

\newpage
% ===========================================================================
% PDF page 242 -- Book page 220
% ===========================================================================
\setcounter{page}{220}

It is hardly necessary to observe, that if
\[
A\xi + B\eta + C\zeta = 0 \qquad (75),
\]
\[
Ax + By + Cz = 0 \qquad (76),
\]
represent the same great circle,
\[
A : B : C = A + \nu B + \mu C : \nu A + B + \lambda C : \mu A + \lambda B + C \qquad (77),
\]
\[
A : B : C = \mathfrak{a}A + \mathfrak{h}B + \mathfrak{g}C : \mathfrak{h}A + \mathfrak{b}B + \mathfrak{f}C : \mathfrak{g}A + \mathfrak{f}B + \mathfrak{c}C \qquad (78).
\]

\subsection*{Inclination of two great Circles.}

Let the equations of these be
\[
A\xi + B\eta + C\zeta = 0 \qquad (79),
\]
\[
\text{or } Ax + By + Cz = 0 \qquad (80),
\]
\[
A'\xi + B'\eta + C'\zeta = 0 \qquad (81),
\]
\[
\text{or } A'x + B'y + C'z = 0 \qquad (82),
\]
and let $e, \mathbf{e}$, have the same values as above, and $e', \mathbf{e}'$, corresponding ones. To obtain the inclination of the two circles, we have only, in the formul\ae{} given above for the distance of two points, to change $p, p', \omega, \omega'$, into $e, e', \mathbf{e}, \mathbf{e}'$.

The distance of a point from a given circle may be found with equal facility; for this is evidently the complement of the distance of the point from the pole of the circle. In like manner we may find the condition that two great circles intersect at right angles, \&c.

There are evidently a whole class of formul\ae, not by any means peculiar to the present system of coordinates, such as
\[
Ax + By + Cz - s(A'x + B'y + C'z) = 0 \qquad (83),
\]
for the equation of a great circle subjected to pass through the points of intersection of
\[
Ax + By + Cz = 0,\ A'x + B'y + C'z = 0.
\]
Again,
\[
\left| \begin{array}{ccc}
x & y & z\\
a & b & c\\
a' & b' & c'
\end{array} \right| = 0 \qquad (84),
\]
for the equation of the great circle which passes through the points given by the sine systems $a : b : c$ and $a' : b' : c'$, \&c., and which are obtained so easily that it is not worth while writing down any more of them.

\subsection*{Transformation of Coordinates.}

Let $X_1, Y_1, Z_1$, be the new points of reference, and suppose $X_1$ is given by the conjoint systems $e = a, b, c$, $\mathbf{e} = \alpha, \beta, \gamma$; and similarly $Y_1, Z_1$, by the analogous systems $e', \mathbf{e}'$; $e'', \mathbf{e}''$.

\newpage
% ===========================================================================
% PDF page 243 -- Book page 221
% ===========================================================================
\setcounter{page}{221}

Suppose, as before, $P$ is given by one of the systems $\omega, p$; and let $\omega_1, p_1$ be the new systems which determine the position of $P$ with reference to $X_1, Y_1, Z_1$.

In the first place, $\lambda_1, \mu_1, \nu_1$, are given by the formul\ae
\[
\lambda_1 = \frac{W(e', e'', q)}{\sqrt{\{W(e', e', q)\,W(e'', e'', q)\}}} = \frac{W(\mathbf{e}', \mathbf{e}'', \mathbf{q})}{\sqrt{\{W(\mathbf{e}', \mathbf{e}', \mathbf{q})\,W(\mathbf{e}'', \mathbf{e}'', \mathbf{q})\}}} \qquad (85),
\]
\[
\mu_1 = \frac{W(e'', e, q)}{\sqrt{\{W(e'', e'', q)\,W(e, e, q)\}}} = \frac{W(\mathbf{e}'', \mathbf{e}, \mathbf{q})}{\sqrt{\{W(\mathbf{e}'', \mathbf{e}'', \mathbf{q})\,W(\mathbf{e}, \mathbf{e}, \mathbf{q})\}}},
\]
\[
\nu_1 = \frac{W(e, e', q)}{\sqrt{\{W(e, e, q)\,W(e', e', q)\}}} = \frac{W(\mathbf{e}, \mathbf{e}', \mathbf{q})}{\sqrt{\{W(\mathbf{e}, \mathbf{e}, \mathbf{q})\,W(\mathbf{e}', \mathbf{e}', \mathbf{q})\}}}.
\]
The system $\omega_1$ is evidently given immediately by
\[
\xi_1 : \eta_1 : \zeta_1 = \frac{W(e, p, q)}{\sqrt{\{W(e, e, q)\}}} : \frac{W(e', p, q)}{\sqrt{\{W(e', e', q)\}}} : \frac{W(e'', p, q)}{\sqrt{\{W(e'', e'', q)\}}} \qquad (86),
\]
\[
= \frac{W(\mathbf{e}, \omega, \mathbf{q})}{\sqrt{\{W(\mathbf{e}, \mathbf{e}, \mathbf{q})\}}} : \frac{W(\mathbf{e}', \omega, \mathbf{q})}{\sqrt{\{W(\mathbf{e}', \mathbf{e}', \mathbf{q})\}}} : \frac{W(\mathbf{e}'', \omega, \mathbf{q})}{\sqrt{\{W(\mathbf{e}'', \mathbf{e}'', \mathbf{q})\}}} \qquad (87),
\]
and from these we may obtain the system $p_1$, by means of the formul\ae
\[
x_1 : y_1 : z_1 = \mathfrak{a}_1\xi_1 + \mathfrak{h}_1\eta_1 + \mathfrak{g}_1\zeta_1 : \mathfrak{h}_1\xi_1 + \mathfrak{b}_1\eta_1 + \mathfrak{f}_1\zeta_1 : \mathfrak{g}_1\xi_1 + \mathfrak{f}_1\eta_1 + \mathfrak{c}_1\zeta_1 \qquad (88).
\]
This requires some further development however. We must in the first place form the system $\mathfrak{a}_1, \mathfrak{b}_1, \mathfrak{c}_1, \mathfrak{f}_1, \mathfrak{g}_1, \mathfrak{h}_1$: this is done immediately from the formul\ae{} of Sect.~2, and we have
\[
\mathfrak{a}_1 = \frac{W(\varpi', \varpi', q)}{W(e', e', q)\,W(e'', e'', q)} = \frac{kW(\bar\varpi', \bar\varpi', \mathbf{q})}{W(\mathbf{e}', \mathbf{e}', \mathbf{q})\,W(\mathbf{e}'', \mathbf{e}'', \mathbf{q})} \qquad (89),
\]
\[
\mathfrak{b}_1 = \frac{W(\varpi'', \varpi'', q)}{W(e'', e'', q)\,W(e, e, q)} = \frac{kW(\bar\varpi'', \bar\varpi'', \mathbf{q})}{W(\mathbf{e}'', \mathbf{e}'', \mathbf{q})\,W(\mathbf{e}, \mathbf{e}, \mathbf{q})},
\]
\[
\mathfrak{c}_1 = \frac{W(\varpi, \varpi, q)}{W(e, e, q)\,W(e', e', q)} = \frac{kW(\bar\varpi, \bar\varpi, \mathbf{q})}{W(\mathbf{e}, \mathbf{e}, \mathbf{q})\,W(\mathbf{e}', \mathbf{e}', \mathbf{q})},
\]
\[
\mathfrak{f}_1 = \frac{W(\varpi'', \varpi, q)}{W(e, e, q)\sqrt{\{W(e', e', q)\,W(e'', e'', q)\}}} = \frac{kW(\bar\varpi'', \bar\varpi, \mathbf{q})}{W(\mathbf{e}, \mathbf{e}, \mathbf{q})\sqrt{\{W(\mathbf{e}', \mathbf{e}', \mathbf{q})\,W(\mathbf{e}'', \mathbf{e}'', \mathbf{q})\}}},
\]
\[
\mathfrak{g}_1 = \frac{W(\varpi, \varpi'', q)}{W(e', e', q)\sqrt{\{W(e'', e'', q)\,W(e, e, q)\}}} = \frac{kW(\bar\varpi, \bar\varpi'', \mathbf{q})}{W(\mathbf{e}', \mathbf{e}', \mathbf{q})\sqrt{\{W(\mathbf{e}'', \mathbf{e}'', \mathbf{q})\,W(\mathbf{e}, \mathbf{e}, \mathbf{q})\}}},
\]
\[
\mathfrak{h}_1 = \frac{W(\varpi', \varpi, q)}{W(e'', e'', q)\sqrt{\{W(e, e, q)\,W(e', e', q)\}}} = \frac{kW(\bar\varpi', \bar\varpi, \mathbf{q})}{W(\mathbf{e}'', \mathbf{e}'', \mathbf{q})\sqrt{\{W(\mathbf{e}, \mathbf{e}, \mathbf{q})\,W(\mathbf{e}', \mathbf{e}', \mathbf{q})\}}}.
\]

\newpage
% ===========================================================================
% PDF page 244 -- Book page 222
% ===========================================================================
\setcounter{page}{222}

\[
x_1 : y_1 : z_1 = \sqrt{\{W(e, e, q)\}} \times
\]
\[
\{W(e, p, q)\,W(\varpi'', \varpi', q) + W(e', p, q)\,W(\varpi'', e'', q) + W(e'', p, q)\,W(\varpi', e', q)\}
\]
\[
: \sqrt{\{W(e', e', q)\}} \times
\]
\[
\{W(e, p, q)\,W(\varpi', \varpi'', q) + W(e', p, q)\,W(\varpi'', \varpi'', q) + W(e'', p, q)\,W(\bar\varpi', e'', q)\}
\]
\[
: \sqrt{\{W(e'', e'', q)\}} \times
\]
\[
\{W(e, p, q)\,W(\varpi, \varpi', q) + W(e', p, q)\,W(\bar\varpi, e', q) + W(e'', p, q)\,W(\bar\varpi', \bar\varpi'', q)\} \qquad (90);
\]
these may be reduced to the very simple form
\[
x_1 : y_1 : z_1 = \sqrt{\{W(e, e, q)\}}\,W(\varpi', \varpi'', q) \qquad (91),
\]
\[
: \sqrt{\{W(e', e', q)\}}\,W(\varpi'', \varpi, q),
\]
\[
: \sqrt{\{W(e'', e'', q)\}}\,W(\varpi, \varpi', q),
\]
and in like manner we obtain
\[
x_1 : y_1 : z_1 = \sqrt{\{W(\mathbf{e}, \mathbf{e}, \mathbf{q})\}}\,W(\bar\varpi', \bar\varpi'', \mathbf{q}) \qquad (92),
\]
\[
: \sqrt{\{W(\mathbf{e}', \mathbf{e}', \mathbf{q})\}}\,W(\bar\varpi'', \bar\varpi, \mathbf{q}),
\]
\[
: \sqrt{\{W(\mathbf{e}'', \mathbf{e}'', \mathbf{q})\}}\,W(\bar\varpi, \bar\varpi', \mathbf{q}).
\]

It will be as well to indicate the steps of this reduction. Consider the quantity in $\{\ \}$ in the first line of the equation which gives the ratios $x_1 : y_1 : z_1$, and suppose for a moment $\overline{e'e''} = l, m, n$, \&c.: then, selecting the portion of the expression which is multiplied by $\mathfrak{a}$, this is
\[
= \mathfrak{a}l\,\{l(a\xi + b\eta + c\zeta) + l'(a'\xi + b'\eta + c'\zeta) + l''(a''\xi + b''\eta + c''\zeta)\},
\]
or, since
\[
la + l'a' + l''a'' = \overline{ee'e''},\quad lb + l'b' + l''b'' = 0,\quad lc + l'c' + l''c'' = 0,
\]
this reduces itself to $\overline{ee'e''}\cdot l\mathfrak{a}\xi$, which is a term of
\[
\overline{ee'e''}\,W(\overline{e'e''}, \omega, q);
\]
and by comparing the remaining terms in the same manner, it would be seen that the whole reduces itself to
\[
\overline{ee'e''}\,W(\overline{e'e''}, \omega, q);
\]
whence the formul\ae{} in question.

The formul\ae{} (86), (87), (91), (92), completely resolve the problem of the transformation of coordinates; they determine respectively $p_1$ from $p$ or $\omega$, $\omega_1$ from $p$ or $\omega$.

To complete the present part of the subject we may add the following formul\ae.

Suppose
\begin{align*}
x_1 : y_1 : z_1 &= ax_1 + a'y_1 + a''z_1 \qquad (93),\\
&: bx_1 + b'y_1 + b''z_1,\\
&: cx_1 + c'y_1 + c''z_1,
\end{align*}

\newpage
% ===========================================================================
% PDF page 245 -- Book page 223
% ===========================================================================
\setcounter{page}{223}

which we see from the preceding formul\ae{} is the form of the relation between the systems $p_1$ and $p$. And suppose, as before, $\lambda_1, \mu_1, \nu_1$ are the cosines of the distances of the new points of reference $X_1, Y_1, Z_1$.

We can immediately determine the relations that must exist between these coefficients, in order that they may actually denote such a transformation. For this purpose write
\begin{align*}
a, b, c &= j \qquad (94),\\
a', b', c' &= j',\\
a'', b'', c'' &= j''.
\end{align*}
Then the distance between the point $P$ and any other point $P'$ is given by the formula
\[
\cos\delta = \frac{W(p, p', \Theta)}{\sqrt{\{W(p, p, \Theta)\,W(p', p', \Theta)\}}} \qquad (95),
\]
where
\[
\Theta = W(j, j, q), W(j', j', q), W(j'', j'', q), W(j', j'', q), W(j'', j, q), W(j, j', q) \qquad (96).
\]

But we must evidently have
\[
\cos\delta = \frac{W(p, p', q)}{\sqrt{\{W(p, p, q)\,W(p', p', q)\}}} \qquad (97),
\]
or the quantities $\Theta$ must be proportional to the quantities $q$, i.e.
\[
W(j, j, q) : W(j', j', q) : W(j'', j'', q) : W(j', j'', q) : W(j'', j, q) : W(j, j', q)
\]
\[
= 1 : 1 : 1 : \lambda_1 : \mu_1 : \nu_1 \qquad (98).
\]

And in precisely the same manner, if instead of $x, y, z, x_1, y_1, z_1$, in the above formul\ae, we had had $\xi, \eta, \zeta : \xi_1, \eta_1, \zeta_1$, the result would have been
\[
W(j, j, \mathbf{q}) : W(j', j', \mathbf{q}) : W(j'', j'', \mathbf{q}) : W(j', j'', \mathbf{q}) : W(j'', j, \mathbf{q}) : W(j, j', \mathbf{q})
\]
\[
= \mathfrak{a} : \mathfrak{b} : \mathfrak{c} : \mathfrak{f} : \mathfrak{g} : \mathfrak{h} \qquad (99).
\]

It is hardly necessary to remark, that throughout the preceding formul\ae{} an expression, such as $W(p, p', q)$, is proportional to either of the quantities
\[
x\xi' + y\eta' + z\zeta'\ \text{ or }\ x'\xi + y'\eta + z'\zeta,
\]
and may be changed into one of these multiplied by an arbitrary constant, which may be always made to disappear by a corresponding change in another quantity of the same form: thus, for instance,
\[
\frac{W(p, p', q)}{W(p, p, q)} = \frac{x\xi' + y\eta' + z\zeta'}{x\xi + y\eta + z\zeta} \qquad (100);
\]
but these forms being unsymmetrical, it is better in general not to introduce them.

All the preceding expressions simplify exceedingly, reducing themselves in fact to the ordinary formul\ae{} for the transformation of rectangular coordinates in Geometry of three dimensions, for the case where the triangle $XYZ$ has its sides and angles right angles. As this presents no difficulty, I shall not enter upon it at present.

\newpage
% ===========================================================================
% PDF page 246 -- Book page 224
% Paper [33] -- ON THE REDUCTION OF du/sqrt(U), WHEN U IS A FUNCTION OF THE FOURTH ORDER
% ===========================================================================
\section*{33.}
\section*{ON THE REDUCTION OF $\dfrac{du}{\sqrt{U}}$, WHEN $U$ IS A FUNCTION OF THE FOURTH ORDER.}
\addcontentsline{toc}{section}{33. On the Reduction of $du/\sqrt{U}$, when $U$ is a Function of the Fourth Order}

\textit{[From the Cambridge and Dublin Mathematical Journal, vol.~i. (1846), pp.~70--73.]}

\bigskip
\setcounter{page}{224}

It is well known that the transformation of this differential expression into a similar one, in which the function in the denominator contains only even powers of the corresponding variable, is the first step in the process of reducing $\int \frac{du}{\sqrt{U}}$ to elliptic integrals. And, accordingly, the different modes of effecting this have been examined, more or less, by most of those who have written on the subject. The simplest supposition, that adopted by Legendre, and likewise discussed in some detail by Gudermann, is that $u$ is a fraction, the numerator and denominator of which are linear functions of the new variable. But the theory of this transformation admits of being developed further than it has yet been done, as regards the equation which determines the modulus of the elliptic function. This may be effected most easily as follows.

Suppose
\[
U = a + 4bu + 6cu^2 + 4du^3 + eu^4,
\]
Also let
\[
P = ax^4 + 4bx^3 y + 6cx^2 y^2 + 4dxy^3 + ey^4.
\]
\[
P' = a'x'^4 + 4b'x'^3 y' + 6c'x'^2 y'^2 + 4d'x'y'^3 + e'y'^4
\]
be what $P$ becomes after writing
\[
x = \lambda x' + \mu y',
\]
\[
y = \lambda_1 x' + \mu_1 y':
\]
and let
\[
U' = a' + 4b'u' + 6c'u'^2 + 4d'u'^3 + e'u'^4.
\]

\newpage
% ===========================================================================
% PDF page 247 -- Book page 225
% ===========================================================================
\setcounter{page}{225}

Suppose, moreover,
\begin{align*}
k &= \lambda\mu_1 - \lambda_1\mu,\\
I &= ae - 4bd + 3c^2,\\
I' &= a'e' - 4b'd' + 3c'^2,\\
J &= ace - ad^2 - b^2 e - c^3 + 2bcd,\\
J' &= a'c'e' - a'd'^2 - b'^2 e' - c'^3 + 2b'c'd';
\end{align*}
we have evidently
\[
xdy - ydx = k(x'dy' - y'dx'),
\]
\[
\frac{xdy - ydx}{P^{1/4}} = \frac{x'dy' - y'dx'}{P'^{1/4}}.
\]
Hence writing
\[
u = \frac{x}{y},\quad u' = \frac{x'}{y'},
\]
and therefore
\[
\frac{xdy - ydx}{P^{1/4}} = \frac{du}{U^{1/4}},\quad \frac{x'dy' - y'dx'}{P'^{1/4}} = \frac{du'}{U'^{1/4}},
\]
we obtain
\[
\frac{du}{U^{1/4}} = k\,\frac{du'}{U'^{1/4}},
\]
the equation between $u$ and $u'$ being
\[
u = \frac{\lambda + \mu u'}{\lambda_1 + \mu_1 u'}.
\]

Next, to determine the relations between the coefficients of $U$ and $U'$. Since $P, P'$ are obtained from each other by linear transformations (Math.~Journal, vol.~IV.~p.~208), [13, p.~94], we have between the coefficients of these functions and of the transforming equations, the relations
\[
I' = k^4 I,
\]
\[
J' = k^6 J.
\]
Suppose now
\[
\frac{J'^2}{I'^3} = \frac{J^2}{I^3}.
\]
or
\[
U' = a'(1 + pu'^2)(1 + qu'^2),
\]
whence also
\[
b' = 0,\quad d' = 0,\quad 6c' = a'(p + q),\quad e' = a'pq;
\]
\[
I' = \tfrac{1}{12} a'^2 (p^2 + q^2 + 14 pq),
\]
\[
J' = \tfrac{1}{216} a'^3 (p + q)(36 pq - p^2 - q^2);
\]

\newpage
% ===========================================================================
% PDF page 248 -- Book page 226
% ===========================================================================
\setcounter{page}{226}

therefore
\[
p^2 + q^2 + 14 pq = 12 \frac{I'}{a'^2},
\]
\[
(p + q)(36 pq - p^2 - q^2) = 216 \frac{J'}{a'^3};
\]
whence also
\[
\frac{(p+q)^2 (36pq - p^2 - q^2)^2}{(p^2 + q^2 + 14 pq)^3} = \frac{27 J'^2}{I'^3},
\]
\[
\frac{108 pq\,(p-q)^2}{(p^2 + q^2 + 14 pq)^2} = 1 - \frac{27 J'^2}{I'^3},
\]
which determines the relation between $p$ and $q$. Also
\[
\frac{k}{\sqrt{a}} = \sqrt[4]{\left( \frac{p^2 + q^2 + 14 pq}{12 I}\right)},
\]
so that
\[
\frac{du}{\sqrt{U}} = \left(\frac{p^2 + q^2 + 14 pq}{12 I}\right)^{\!1/4} \cdot \frac{du'}{\sqrt{\{(1 + pu'^2)(1 + qu'^2)\}}}.
\]
If in particular $p = -1$, writing also $-q$ for $q$,
\[
\frac{du}{\sqrt{U}} = \left(\frac{q^2 + 14q + 1}{12}\right)^{\!1/4} \cdot \frac{du'}{\sqrt{\{(1 - u'^2)(1 - qu'^2)\}}},
\]
where
\[
\frac{108 q\,(1 - q)^2}{(q^2 + 14q + 1)^3} = \frac{27 J^2}{I^3}.
\]
Suppose, for shortness,
\[
M = \frac{1}{(1 - \frac{27 J^2}{I^3})},\quad \text{or}\ \frac{27 J^2}{I^3} = 1 - \frac{1}{M};
\]
then
\[
(q^2 + 14 q + 1)^3 - 108 M q\,(q - 1)^2 = 0,
\]
i.e.
\[
\left(q + 1 + \frac{14 q}{q + 1}\right)^{\!3} - 108\,M\,\frac{q\,(q - 1)^2}{(q + 1)^3} = 0.
\]
Let
\[
\frac{q^2 - q + 1}{q + 1} = \theta,
\]
then
\[
\theta^3 - M\,(\theta + 1)^2 = 0,
\]
which determines $\theta$. And then
\[
q = \tfrac12 + \theta + \theta^{\frac12}\sqrt{(3 + \theta)}.
\]

\newpage
% ===========================================================================
% PDF page 249 -- Book page 227
% ===========================================================================
\setcounter{page}{227}

Suppose $q = \alpha$ is one of the values of $q$; the equation becomes
\[
\frac{(q^2 + 14 q + 1)^3}{q(q - 1)^2} = \frac{(\alpha^2 + 14 \alpha + 1)^3}{\alpha(\alpha - 1)^2} = \frac{(\beta^8 + 14\beta^4 + 1)^3}{\beta^4(\beta^4 - 1)^4}.
\]
Now if
\[
q = \beta^4,
\]
then
\[
q^2 + 14 q + 1 = \beta^8 + 14 \beta^4 + 1,
\]
which values satisfy the above equation; hence also, identically,
\[
\frac{(q^2 + 14 q + 1)^3}{q(q - 1)^2} = \frac{(\alpha^2 + 14 \alpha + 1)^3}{\alpha(\alpha - 1)^2},
\]
or the values of $q$ take the form
\[
q = \alpha,\ \beta^4.
\]
(Comp.~Abel.~\OE{}uvres, tom.~i. p.~310 [Ed.~2, p.~459].)

The equation
\[
\theta^3 - M\theta^2 - 2 M\theta - M = 0
\]
has its three real roots if $27 - 4 M$ is negative, and only a single real root if $27 - 4M$ is positive. Writing the equation under the form
\[
(\theta + 3)^3 - 9(\theta + 3)^2 + (27 - M)(\theta + 3) - (27 - 4M) = 0,
\]
the positive roots are both greater than 3. Writing it under the form
\[
(\theta - 1)^3 + 3(\theta - 1)^2 + (3 - M)(\theta - 1) + 1 = 0,\quad (3 - M\ \text{is negative})
\]
the positive roots are both greater than 1. Hence, in this case, $q$ has four positive values. In the former case, $\frac{J^2}{I^3} = \frac{1}{27}\,\frac{M-1}{M}$ is positive, and the function $U$ has either four imaginary factors or four real ones. In the second case, $\frac{J^2}{I^3}$ is negative, or the function $U$ has two real and two imaginary.

\newpage
% ===========================================================================
% PDF page 250 -- Book page 228
% Paper [34] -- NOTE ON THE MAXIMA AND MINIMA OF FUNCTIONS OF THREE VARIABLES.
% ===========================================================================
\section*{34.}
\section*{NOTE ON THE MAXIMA AND MINIMA OF FUNCTIONS OF THREE VARIABLES.}
\addcontentsline{toc}{section}{34. Note on the Maxima and Minima of Functions of Three Variables}

\textit{[From the Cambridge and Dublin Mathematical Journal, vol.~i. (1846), pp.~74, 75.]}

\bigskip
\setcounter{page}{228}

If $A, B, C, F, G, H$, be any real quantities, such that
\[
-BC + CA + AB - F^2 - G^2 - H^2,
\]
and
\[
(A + B + C)(ABC - AF^2 - BG^2 - CH^2 + 2FGH)
\]
are positive; the six quantities
\[
BC - F^2,\ CA - G^2,\ AB - H^2,\ AK,\ BK,\ CK,
\]
(where $K = ABC - AF^2 - BG^2 - CH^2 + 2FGH$) are all of them positive. It is unnecessary to point out the connection of this property with the theory of maxima and minima.

To demonstrate this, writing as usual
\begin{align*}
BC - F^2 &= A',& GH - AF &= F',\\
CA - G^2 &= B',& HF - BG &= G',\\
AB - H^2 &= C',& FG - CH &= H',
\end{align*}
and $K$ as above: then if $A'', B'', C'', F'', G'', H'', K'$ be formed from $A', B', C', F', G', H'$, as these and $K$ are from $A, B, C, F, G, H$, we have the well-known formul\ae
\begin{align*}
A'' &= KA,& F'' &= KF,& K' &= K^2,\\
B'' &= KB,& G'' &= KG,&\\
C'' &= KC,& H'' &= KH.
\end{align*}

\end{document}
